\documentclass[../main]{subfiles}
\begin{document}
\chapter{Motor asincrónico}
\begin{center}
\includegraphics[scale=0.4]{images/04/motor.jpg}
\end{center}

\info{Motor Trifásico}{Deslizamiento,principio de funcionamiento, arranque.}
\actividad{Trabajo práctico de investigación:}
\begin{enumerate}
    \item Explique que es el resbalamiento o deslizamiento en un motor asincronico.
    \item Investigar la secuencia paso a paso del campo magnético giratorio en un motor trifásico.
    \item Explicar la curva característica de un motor trifásico asincrónico. (curva par-velocidad) 
    \item Explicar el rendimiento de un motor y su influencia sobre el par.
    \item Explicar como se calcula el par de un motor sus unidades, además:
    \begin{itemize}
        \item Par máximo.
        \item Par nominal.
        \item Par de arranque.
    \end{itemize}
    \item Explicar la forma de conexión de devanados en un motor trifásico y las formas de conexión en la placa del motor.
    \item Explicar los siguientes sistemas de arranque, realizar esquema de mando y potencia en caso de ser necesario.
    \begin{itemize}
        \item Arranque Directo.
        \item Arranque estrella-triángulo.
        \item Arranque por devanados separados.
        \item Arranque por variador de frecuencia.
        \item Arranque suave.
    \end{itemize}
    \item Investigar información de internet de al menos 3 motores trifásicos, obteniendo, resbalamiento, rendimiento,corrientes nominal, par máximo, par nominal, par de arranque, corriente de arranque, y comparar en una tabla.
    \item Investigación: se desea elevar una carga de $P=1000Kg$, a una velocidad lineal y constante de $v_0=10m/s$,se pueden usar mecanismos de reducción y poleas ¿Qué motor debería comprar?
    \item ¿Si duplico el peso, a que velocidad podrá ascender la carga?
\end{enumerate}
\end{document}
